\scanpage{75}
% archive foliation: 74
\uline{Catégories} \uline{hypopolynomiales}.

Une \uline{catégorie polynomiale} est une catégorie $\unsure{\mathcal P|P}$, munie d'une sous-catégorie (pas pleine en général) $\unsure{\mathcal P|P}_1$, avec les axiomes suivants
\margin{Notation $\unsure{\mathcal P|P}(X,Y)$ pour $\operatorname{Hom}_\unsure{\mathcal P|P}(X,Y)$.}

a) $\unsure{\mathcal P|P}$ admet des produits finis, \deleted{et l'objet final.}

b) $\operatorname{Ob}\unsure{\mathcal P|P}_1=\operatorname{Ob}\unsure{\mathcal P|P}$, \deleted{\unsure{le produit et pour deux objets} dans $\unsure{\mathcal P|P}_1$, leur produit dans $\unsure{\mathcal P|P}$ est un produit dans $\unsure{\mathcal P|P}_1$ (i.e. si $X,Y,Z\in\operatorname{Ob}\unsure{\mathcal P|P}_1$ \unsure{une} \inserted{\unsure{$X\times Y\in\operatorname{Ob}\unsure{\mathcal P|P}_1$}} \unsure{flèche} $Z\longrightarrow X\times Y$ de $\unsure{\mathcal P|P}$ est de $\unsure{\mathcal P|P}_1$ \inserted{(i.e. si $X,Y,Z\in\operatorname{Ob}\unsure{\mathcal P|P}_1$ une $\unsure{\mathcal P|P}$-flèche)} \unsure{$u=(u_1,u_2):X\longrightarrow Y\times Z$ est dans $\unsure{\mathcal P|P}_1$ ssi}}
(donc $\unsure{\mathcal P|P}_1$ est décrite par un ens $\operatorname{Fl}_1\subset\operatorname{Fl}(\unsure{\mathcal P|P})=\operatorname{Fl}$ contenant les identités, et stable par composition.) On veut de plus que \deleted{$\operatorname{Fl}_1$ \unsure{contienne les iso, et que}} pour $X,Y\in\operatorname{Ob}\unsure{\mathcal P|P}$, leur produit $X\times Y$ soit un produit de $\unsure{\mathcal P|P}_1$ (i.e. une flèche $u=(u_1,u_2):Z\longrightarrow X\times Y$ est ``de degré $1$'' ssi ses composantes $X,Y$ le sont), \inserted{enfin si $0$ est l'objet final, tt $X\longrightarrow0$ est de degré $1$ (donc $0$ objet final de $\unsure{\mathcal P|P}_1$)}.

c) $\unsure{\mathcal P|P}_1$ est une catégorie additive.

Cela implique que \deleted{pour tt objet l'on a un hom de monoïdes
\diagram{\begin{tikzcd}
\mathbf Z \arrow[r] & \operatorname{End}(\operatorname{id}\unsure{\mathcal P|P}_1) & \operatorname{End}(\operatorname{id}\unsure{\mathcal P|P}) \arrow[l]
\end{tikzcd}}
\unsure{qui pour tt} $\unsure{n}\in\unsure{\mathbf Z}$, $X,Y\in\operatorname{Ob}\unsure{\mathcal P|P}_1$, \unsure{$n\cdot X$} \inserted{i.e. pour tt $X$}}
les $\operatorname{Hom}_\unsure{\mathcal P|P}(X,Y)$ sont des groupes abéliens (car les objets \deleted{$Y$} \inserted{de $\unsure{\mathcal P|P}$} sont des objets groupes abéliens de $\unsure{\mathcal P|P}$), dépendant fonctoriellt de $X$, mais pas de $Y$ (i.e. on a
\[
\boxed{(f+g)\circ h=f\unsure{\circ}h+g\circ h}
\]
mais \deleted{\unsure{on a en général}} \inserted{on n'a pas} \unsure{la formule}
\[
f\circ(g+h)\ne f\circ g+f\circ h\qquad !
\]
Il en résulte aussi que $\mathbf Z^\times$ opère sur $\unsure{\mathcal P|P}$ i.e. on a un homomorphisme multiplicatif
\diagram{\begin{tikzcd}
\mathbf Z^\times \arrow[r] & \operatorname{End}(\unsure{\mathcal P|P}_1) & \operatorname{End}(\unsure{\mathcal P|P}) \arrow[l,no head]
\end{tikzcd}}
qui à $n\in\mathbf Z$ et $X\in\operatorname{Ob}\unsure{\mathcal P|P}_1=\operatorname{Ob}\unsure{\mathcal P|P}$ associe $n\operatorname{id}_X$.
Ainsi $\mathbf Z$ opère dans \deleted{\unsure{$\operatorname{Hom}$}} $\unsure{\mathcal P|P}(X,Y)$ par
\[
\mu_n(f)=f\circ(n\operatorname{id}_Y)
\]

